\documentclass[11pt,a4paper]{article}

\usepackage[utf8]{inputenc}
\usepackage[T1]{fontenc}
\usepackage[english]{babel}
\usepackage{amsmath,amssymb}
\usepackage[top=1.2cm,bottom=1.5cm,left=2cm,right=2cm]{geometry}
\usepackage{graphicx}
\usepackage{booktabs}
\usepackage{hyperref}
\hypersetup{colorlinks=true,linkcolor=black,citecolor=blue,urlcolor=blue}

\title{Section E --- Threshold Choice and the No-Drift Control}
\author{Alexandre Boyer}
\date{}

\begin{document}
\maketitle

\noindent This section collects results that belong to no single question of the subject
and that would distort the one they were attached to. Two are established here. First,
detecting a drift and winning the race against the forest are different things, and they
separate as soon as the threshold rises. Second, the forest replaces trees at a steady rate
when nothing changes at all, which disqualifies the date of the first replacement as a
signal on its own. Everything below is recomputed offline from the recorded error
trajectories by \texttt{resultats\_R2/scripts/chiffres\_QE.py}; no campaign is rerun.

\paragraph{Control.} $S_t$ does not depend on $\lambda$, so a single offline pass over each
trajectory yields the first crossing for every threshold. The sweep must reproduce the
campaign exactly at the three thresholds of the subject, and it does: $1\,906$, $778$ and
$1$ alarms at $\lambda = 8$, $25$ and $50$. Nothing below is readable without this check.

\section{Detecting Is Not Winning}

The campaign of $2\,000$ runs is read at six thresholds. The drift-free arm of A1
($2\,000$ runs, no change point at all, same horizon) gives what the same threshold costs when
nothing happens.

\begin{center}
\begin{tabular}{crrrl}
\toprule
$\lambda$ & alarm fires & wins the race & of uncensored & false alarms, no drift \\
\midrule
$8$  & $95.3\%$ & $\mathbf{91.0\%}$ & $95.4\%$ & $83/2\,000 = 4.15\%$ $[3.36;5.12]$ \\
$12$ & $94.3\%$ & $78.8\%$ & $83.5\%$ & $11/2\,000 = 0.55\%$ $[0.31;0.98]$ \\
$15$ & $92.9\%$ & $\mathbf{36.3\%}$ & $39.1\%$ & $3/2\,000 = 0.15\%$ $[0.05;0.44]$ \\
$20$ & $66.2\%$ & $8.1\%$ & $12.2\%$ & $0/2\,000 = 0.00\%$ $[0.00;0.19]$ \\
$25$ & $38.9\%$ & $2.5\%$ & $6.6\%$ & $0/2\,000 = 0.00\%$ $[0.00;0.19]$ \\
$50$ & $0.1\%$ & $0.0\%$ & $0.0\%$ & $0/2\,000 = 0.00\%$ $[0.00;0.19]$ \\
\bottomrule
\end{tabular}
\end{center}

\noindent ``Wins the race'' is $\tau_{\mathrm{det}} < \tau_{\mathrm{ARF}}$ over all
$2\,000$ runs, censored runs counted as losses; the next column restricts to the runs that
alarm. Brackets: Wilson $95\%$ interval, $n = 2\,000$ drift-free runs.

\paragraph{The two columns come apart.} At $\lambda = 15$ the alarm still fires in
$92.9\%$ of runs, barely below the $95.3\%$ of $\lambda = 8$, but it arrives first in only
$36.3\%$ of them. Almost two thirds of its alarms are raised after the forest has already
replaced a tree --- the \emph{Zombie alarms} of the manuscript. Reading the detection
column alone would suggest that a moderate threshold escapes the blind spot. It does not.

\paragraph{The paradoxical signature survives at $\lambda = 15$.} Amplitude by amplitude,
the share of runs won by the detector falls as the drift grows, which is the signature the
manuscript reports at $\lambda = 25$.

\begin{center}
\begin{tabular}{crrr}
\toprule
$\Delta e$ & $\lambda = 8$ & $\lambda = 15$ & $\lambda = 20$ \\
\midrule
$0.028$ & $0\%$  & $0\%$  & $0\%$ \\
$0.085$ & $62\%$ & $38\%$ & $20\%$ \\
$0.141$ & $88\%$ & $81\%$ & $49\%$ \\
$0.243$ & $97\%$ & $71\%$ & $21\%$ \\
$0.361$ & $98\%$ & $44\%$ & $3\%$ \\
$0.498$ & $99\%$ & $18\%$ & $1\%$ \\
\bottomrule
\end{tabular}
\end{center}

\noindent At $\lambda = 8$ the rate climbs with the amplitude, from $62\%$ to $99\%$. At
$\lambda = 15$ it peaks at $81\%$ then falls to $18\%$, and at $\lambda = 20$ it collapses
to $1\%$. The transition $\lambda^* \approx 15$ to $25$ announced in the manuscript is
therefore confirmed on the race, not contradicted. Q~D.1, Table~2 gives the same column at
$\lambda = 8$ over all twenty amplitudes.

\paragraph{What the false-alarm column adds.} The manuscript grants that a hyper-reactive
threshold reaches a safe zone but states that it ``sacrifices false-alarm robustness''. On
this bench that cost is measured rather than assumed: $\lambda = 8$ raises an alarm in
$83$ of $2\,000$ runs where nothing changes, $4.15\%$, just under the $5\%$ calibration level usually adopted for
drift detectors. So on the Bernoulli bench the only threshold that wins the race does so at
an ordinary false-alarm cost, and the admissible set is not empty here.

\paragraph{Three reservations.} The false-alarm rate is measured over $2\,000$ steps on
$2\,000$ runs; over a longer deployment such alarms accumulate, and the upper end of the
Wilson interval at $\lambda = 8$ is $5.12\%$. The bench is homoscedastic, so this says
nothing about the \emph{Fundamental Tension} the manuscript establishes on
ARMA--GARCH streams, where bounding false alarms under clustered volatility is what empties
the admissible set. And at the weakest amplitude no threshold wins a single run, which
Q~A establishes without any dependence assumption.

\section{The No-Drift Control}

The same drift-free arm answers a question that needs no threshold at all: what does
$\tau_{\mathrm{ARF}}$ read when there is nothing to repair?

\begin{center}
\begin{tabular}{lr}
\toprule
On $2\,000$ runs with no change point & \\
\midrule
runs where at least one tree is replaced & $2\,000/2\,000$ \\
first replacement, median & $95$ steps (quartiles $39$ and $192$) \\
distinct trees renewed over the horizon, median & $9/10$ \\
runs whose forest is entirely renewed & $575/2\,000 = 28.7\%$ \\
\midrule
median $\tau_{\mathrm{ARF}}$ at $\Delta e = 0.028$ (real drift) & $118$ steps \\
median $\tau_{\mathrm{ARF}}$ over the whole grid & $38$ steps \\
\bottomrule
\end{tabular}
\end{center}

\paragraph{Reading.} Replacement churn under pure noise is the normal regime of the ARF,
not an anomaly: nine trees out of ten are renewed over two thousand steps with no drift
whatsoever. At the weakest amplitude of the grid the first replacement comes after $118$ steps against
$95$ with no drift at all, so at that amplitude it does not distinguish a real change from
its absence. This is the control any candidate replacement for $\tau_{\mathrm{ARF}}$ must
pass before its correlations are worth computing, and it is test A1 of the benchmark.

\section{What This Section Establishes}

Detection and victory separate: a threshold can alarm on almost every drift and still
arrive last on two thirds of them, so a detection rate quoted alone does not describe the
blind spot. On this bench one threshold does win the race, $\lambda = 8$, at a false-alarm
cost of $4.15\%$ per two thousand drift-free steps, which restricts the manuscript's
``no CUSUM threshold escapes it'' to the settings it actually tested here. And the date of
the first replacement fires as readily without drift as with the weakest one, which is a
property of the indicator rather than of the detector.

\end{document}
